\documentclass[10pt,a4paper]{amsart}
\usepackage{amsmath,amssymb,amsthm,graphicx}
% Retained minimal preamble for the Stacks Project chapter sites-cohomology.tex (Cohomology on Sites).
% Verbatim excerpt of the chapter preamble: Xy-pic setup (native preamble lines 14-18) and the native
% theorem declarations with the mathematical macros (native preamble lines 154-217). Copied mechanically;
% nothing added, removed or redefined. The wrapper class and the amsmath/amssymb/amsthm/graphicx base
% packages are supplied by the assembler, so the class/externaldocument part of the native preamble is omitted.
% For commutative diagrams we use Xy-pic
\usepackage[all]{xy}

% We use 2cell for 2-commutative diagrams.
\xyoption{2cell}
% Theorem environments.
%
\theoremstyle{plain}
\newtheorem{theorem}[subsection]{Theorem}
\newtheorem{proposition}[subsection]{Proposition}
\newtheorem{lemma}[subsection]{Lemma}

\theoremstyle{definition}
\newtheorem{definition}[subsection]{Definition}
\newtheorem{example}[subsection]{Example}
\newtheorem{exercise}[subsection]{Exercise}
\newtheorem{situation}[subsection]{Situation}

\theoremstyle{remark}
\newtheorem{remark}[subsection]{Remark}
\newtheorem{remarks}[subsection]{Remarks}

\numberwithin{equation}{subsection}

% Macros
%
\def\lim{\mathop{\mathrm{lim}}\nolimits}
\def\colim{\mathop{\mathrm{colim}}\nolimits}
\def\Spec{\mathop{\mathrm{Spec}}}
\def\Hom{\mathop{\mathrm{Hom}}\nolimits}
\def\Ext{\mathop{\mathrm{Ext}}\nolimits}
\def\SheafHom{\mathop{\mathcal{H}\!\mathit{om}}\nolimits}
\def\SheafExt{\mathop{\mathcal{E}\!\mathit{xt}}\nolimits}
\def\Sch{\mathit{Sch}}
\def\Mor{\mathop{\mathrm{Mor}}\nolimits}
\def\Ob{\mathop{\mathrm{Ob}}\nolimits}
\def\Sh{\mathop{\mathit{Sh}}\nolimits}
\def\NL{\mathop{N\!L}\nolimits}
\def\CH{\mathop{\mathrm{CH}}\nolimits}
\def\proetale{{pro\text{-}\acute{e}tale}}
\def\etale{{\acute{e}tale}}
\def\QCoh{\mathit{QCoh}}
\def\Ker{\mathop{\mathrm{Ker}}}
\def\Im{\mathop{\mathrm{Im}}}
\def\Coker{\mathop{\mathrm{Coker}}}
\def\Coim{\mathop{\mathrm{Coim}}}

% Boxtimes
%
\DeclareMathSymbol{\boxtimes}{\mathbin}{AMSa}{"02}

%
% Macros for moduli stacks/spaces
%
\def\QCohstack{\mathcal{QC}\!\mathit{oh}}
\def\Cohstack{\mathcal{C}\!\mathit{oh}}
\def\Spacesstack{\mathcal{S}\!\mathit{paces}}
\def\Quotfunctor{\mathrm{Quot}}
\def\Hilbfunctor{\mathrm{Hilb}}
\def\Curvesstack{\mathcal{C}\!\mathit{urves}}
\def\Polarizedstack{\mathcal{P}\!\mathit{olarized}}
\def\Complexesstack{\mathcal{C}\!\mathit{omplexes}}
% \Pic is the operator that assigns to X its picard group, usage \Pic(X)
% \Picardstack_{X/B} denotes the Picard stack of X over B
% \Picardfunctor_{X/B} denotes the Picard functor of X over B
\def\Pic{\mathop{\mathrm{Pic}}\nolimits}
\def\Picardstack{\mathcal{P}\!\mathit{ic}}
\def\Picardfunctor{\mathrm{Pic}}
\def\Deformationcategory{\mathcal{D}\!\mathit{ef}}
\title{An ascending chain condition on a graded piece forces the Mittag-Leffler condition for the cohomology of an inverse system of sheaves}
\author{The Stacks Project Authors}
\date{stacks-project:sites-cohomology.tex:local-snapshot-20261011; GNU FDL 1.2 or later}
\begin{document}
\maketitle
\noindent\textit{Editorial note.} All delivered mathematical text is the author's own, copied from the pinned Stacks Project chapter file sites-cohomology.tex; no statement or proof step was rewritten, reordered or removed.\par\medskip
\noindent\textit{Editorial note.} The proof is the author's complete argument, including the graded compatibility diagram and the final stable-image computation; the author's closing sentence is retained verbatim.\par\medskip
\noindent\textit{Editorial note.} The stated lemma carries a footnote spelling out the stable-image form of Mittag-Leffler; the footnote lies inside the counted statement span and is delivered as the author wrote it.\par\medskip
\noindent\textit{Editorial note.} Cross-chapter restatement: the same label occurs in chapter cohomology (topological spaces) and in chapter more-algebra. The conclusion is the same and those statements are special cases of this site statement, so per the lane-s4 rule only this more general occurrence is counted here.\par\medskip
\noindent\textit{Editorial note.} The preamble is the chapter's own retained minimal preamble (sources/stacks-sites-cohomology/retained-preamble.tex), a verbatim excerpt of the native preamble: the Xy-pic setup and the native theorem declarations with the mathematical macros.\par\medskip
\noindent\textit{Allowed standard prerequisites (editorial).}\begin{itemize}
\item Mittag-Leffler condition for an inverse system of modules. Same chapter, the Mittag-Leffler condition as the stabilisation of the images Im(M\_n -> M\_\{n-c\}). Supplies the conclusion, in the equivalent stable-image form stated in the lemma's footnote.
\item Long exact cohomology sequence of a short exact sequence of sheaves. Same chapter, the boundary map delta\_n : M\_n -> N\_n attached to 0 -> I\textasciicircum{}n F\_\{n+1\} -> F\_\{n+1\} -> F\_n -> 0. Supplies the boundary maps whose images are shown to form a finitely generated graded module.
\item Graded module structure over the Rees-type ring of the powers of I. Same chapter, the graded module structure of the direct sum of the N\_n over the direct sum of the I\textasciicircum{}n/I\textasciicircum{}\{n+1\}. Supplies the ascending chain condition hypothesis, which is applied to the graded submodule of the images of the delta\_n.
\item Commutative diagram relating the boundary maps for n and n + m. Same chapter, the morphism of short exact sequences induced by multiplication by an element f of I\textasciicircum{}m, giving delta\_\{n+m\}(fs) = f delta\_n(s). Supplies the compatibility that makes the direct sum of the images of the delta\_n a graded submodule.
\end{itemize}\medskip
\begin{lemma}
\label{lemma-ML-general}
Let $I$ be an ideal of a ring $A$. Let $\mathcal{C}$ be a site.
Let
$$
\ldots \to \mathcal{F}_3 \to \mathcal{F}_2 \to \mathcal{F}_1
$$
be an inverse system of sheaves of $A$-modules on $\mathcal{C}$
such that $\mathcal{F}_n = \mathcal{F}_{n + 1}/I^n\mathcal{F}_{n + 1}$.
Let $p \geq 0$. Assume
$$
\bigoplus\nolimits_{n \geq 0} H^{p + 1}(\mathcal{C}, I^n\mathcal{F}_{n + 1})
$$
satisfies the ascending chain condition as a graded
$\bigoplus_{n \geq 0} I^n/I^{n + 1}$-module.
Then the inverse system $M_n = H^p(\mathcal{C}, \mathcal{F}_n)$ satisfies the
Mittag-Leffler condition\footnote{In fact, there exists
a $c \geq 0$ such that $\Im(M_n \to M_{n - c})$ is the stable image
for all $n \geq c$.}.
\end{lemma}

\begin{proof}
Set $N_n = H^{p + 1}(\mathcal{C}, I^n\mathcal{F}_{n + 1})$ and let
$\delta_n : M_n \to N_n$ be the boundary map on cohomology
coming from the short exact sequence
$0 \to I^n\mathcal{F}_{n + 1} \to \mathcal{F}_{n + 1} \to \mathcal{F}_n \to 0$.
Then $\bigoplus \Im(\delta_n) \subset \bigoplus N_n$ is a graded submodule.
Namely, if $s \in M_n$ and $f \in I^m$, then we have a commutative diagram
$$
\xymatrix{
0 \ar[r] &
I^n\mathcal{F}_{n + 1} \ar[d]_f \ar[r] &
\mathcal{F}_{n + 1} \ar[d]_f \ar[r] &
\mathcal{F}_n \ar[d]_f \ar[r] & 0 \\
0 \ar[r] &
I^{n + m}\mathcal{F}_{n + m + 1} \ar[r] &
\mathcal{F}_{n + m + 1} \ar[r] &
\mathcal{F}_{n + m} \ar[r] & 0
}
$$
The middle vertical map is given by lifting a local section of
$\mathcal{F}_{n + 1}$ to a section of $\mathcal{F}_{n + m + 1}$
and then multiplying by $f$; similarly for the other vertical arrows.
We conclude that $\delta_{n + m}(fs) = f \delta_n(s)$.
By assumption we can find $s_j \in M_{n_j}$, $j = 1, \ldots, N$
such that $\delta_{n_j}(s_j)$
generate $\bigoplus \Im(\delta_n)$ as a graded module. Let $n > c = \max(n_j)$.
Let $s \in M_n$. Then we can find $f_j \in I^{n - n_j}$ such that
$\delta_n(s) = \sum f_j \delta_{n_j}(s_j)$. We conclude that
$\delta(s - \sum f_j s_j) = 0$, i.e., we can find $s' \in M_{n + 1}$
mapping to $s - \sum f_js_j$ in $M_n$. It follows that
$$
\Im(M_{n + 1} \to M_{n - c}) = \Im(M_n \to M_{n - c})
$$
Namely, the elements $f_js_j$ map to zero in $M_{n - c}$.
This proves the lemma.
\end{proof}

\end{document}
